\section{Relation to prior results and the scope of the theorem}\label{sec:comparison}
The present result concerns a finite-state implementation of a complete experiment, rather than only an approximation to a Bellman function. The distinction does not make nearby belief-aggregation work irrelevant. Grover and Dimitrakakis~\cite{grover}, Theorem~1 and Lemma~1, use depth-dependent discretizations to control planning error and the size of a belief tree. Thus adapting resolution to continuation sensitivity is not by itself a new principle. Here the objects counted are deployed autonomous labels; actual conditional barycenters give a second-order task loss, and a clock converse matches that construction.

Li, Hammar and Bertsekas~\cite{li}, Proposition~5, bound the error of a two-stage feature/belief aggregate model by an oscillation quantity divided by one minus the discount factor. Their subsequent lower-approximation condition is a statement about the aggregate cost function, not a decision-theoretic lower for every $M$-label observer. Our analysis uses the posterior law of the implemented observer itself and a lower from its acquired mass. The two constructions answer different questions; neither is described as a special case of the other without an explicit reduction.

Kara and Y\"uksel~\cite{kara}, Theorem~12, relate finite-window policy error to filter stability, with exponential convergence under stated conditions and a discount restriction. Demirci, Kara and Y\"uksel~\cite{demirci} obtain contraction-based average-cost optimality and approximation results. These are important precedents for using stability to suppress remote information. The finite exact-deadline task here is not an average-cost or infinite-horizon result: the new joint bound includes the necessary $n$ clock states and the remaining $(M-n)$ resolution states, while an exact calibrated Jensen identity supplies the task error. The elementary factor-three instrument coupling is not claimed as a sharper replacement for those filter results.

Lyu, Suen and Zhang~\cite{lyu}, Theorem~1 and Proposition~1, characterize static optimal interval partitions using conditional means and value curvature, including uniqueness under logconcavity. Consequently conditional-mean aggregation and curvature-sensitive partition design are established ideas. The local quantity \eqref{eq:masscurvature} serves a different role: it bounds sequential losses under a machine-dependent law, permits nonsmooth continuation values and boundary approach, and is combined with explicit online state accounting. It is sufficient, not a characterization of all optimal partitions.

The operator/instrument language of Davies and Lewis~\cite{davies}, finite-dimensional Bayesian conditioning, supporting hyperplanes, small-ball quantization and Blackwell--Le Cam comparison~\cite{blackwell,lecam} are standard tools. The proof above supplies the specific statements used, rather than attributing those tools to the present work. The noncommuting inverse \eqref{eq:quantuminverse} and the Gaussian density calculation verify this theorem's experimental hypotheses; they do not turn matrix parameter counting into an acquired-geometry theorem.

\subsection{What is and is not combined}
Theorem~\ref{thm:main} combines a sharp clock converse with a horizon-uniform causal upper on a finite affine predictive class. Theorem~\ref{thm:strata} adds actual singular mass and the two raw-kernel theorems verify the same general hypotheses without reversing logical dependence. The exact identity is valid beyond the density and strong-concavity assumptions; the explicit exponent is not.

Supplement U retains the complete posterior-pooling article, including its strictly adaptive affine-sensor example and finite-precision arguments. Supplement T retains the complete acquired-refinement and serial-continuation article: its fixed-target adaptive acquisition and its noisy intermediate-cut obstructions are not replaced by the present fixed-deadline law. Supplement S retains the complete earlier block-stability, singular-transport and stopped-task proofs. These results have their original hypotheses and independent theorem numbers. Only the specific identity and analytic mechanisms cited in the introduction enter the present proof, and they are proved here in the form needed. Separate ordered-measurement, streaming-memory and finite-action research remains separate.

Several genuine generalization questions remain. Uniform acquisition density is a substantive smoothing condition; the mass--curvature expression suggests weaker hypotheses but does not by itself make them necessary. A general positive experiment need not have a finite affine realization or a Borel predictive quotient without additional measurability assumptions. The block condition is sufficient, not an intrinsic characterization of optimal causal completion, and it is not a theorem for unknown kernel learning. The exact-$n$ task family permits a time-dependent terminal hidden state; it is not automatically a monotone fixed-target sampling law as $n$ changes. Unbounded stopping and infinite time require new estimates. Finally, the exact finite-state construction and its conditional digital certificates do not determine the optimal joint region of programme length, workspace, numerical computation, reset budget and physical time. These are mathematical boundaries of the stated results, not substitutions of a different foundations problem.
